\documentclass[11pt]{article}
\usepackage[T1]{fontenc}
\usepackage[utf8]{inputenc}
\usepackage{lmodern}
\usepackage{amsmath,amssymb,amsthm,mathtools}
\usepackage{array,booktabs}
\usepackage{microtype}
\usepackage[margin=1in]{geometry}
\usepackage[hidelinks]{hyperref}
\usepackage{xcolor}

\newtheorem{theorem}{Theorem}[section]
\newtheorem{proposition}[theorem]{Proposition}
\newcommand{\C}{\mathbb C}
\newcommand{\R}{\mathbb R}
\newcommand{\N}{\mathbb N}
\newcommand{\SigmaL}{\Sigma_L}
\newcommand{\OracleReports}{3,046 reports: 3,019 using only standard axioms and 27 axiom-free; zero nonstandard axioms and zero \texttt{sorry}}
\newcommand{\ModuleHash}{\texttt{c474ba533c25c46f920ce70284980998bcdadeff8c246ec0e38a1240102b3e14}}
\newcommand{\OleanHash}{\texttt{0751d3b593c3f845175cf05e1e863ced46343748796885a76705046d8c3cad0e}}
\newcolumntype{P}[1]{>{\raggedright\arraybackslash}p{#1}}

\title{From Reflection Positivity to a Uniform Reconstructed Gap:\\
A Machine-Checked Transfer Operator for Coupled $\mathbb Z_2$ Slices}
\author{Llu\'is Eriksson}
\date{4 August 2026}

\begin{document}
\maketitle

\begin{abstract}
We close a compositional gap between two previously machine-checked finite-volume
constructions.  The first reconstructs a transfer operator from the positive
site form forced by reflection positivity.  The second proves a
volume-uniform Dobrushin spectral gap for a concrete tilted kernel.  We define
the square-root dressing of real boundary vectors, prove that it is an
isometry for the reconstructed site form, identify the forced transfer
operator with the symmetric weighted kernel, transport this identity through
all finite iterates and Perron normalisation, and finally compose it with the
existing Dobrushin theorem.  Under
\[
  2\tanh |\beta|+2\tanh |\gamma|\leq \alpha<1,
\]
one positive exponent works for every spatial extent.  The six interface
theorems are checked by Lean~4 against Mathlib.  The contribution is the exact
machine-checked composition and its binder order, not a new analytic estimate.
\end{abstract}

\section{Scope and claim discipline}

Osterwalder--Schrader (OS) reflection positivity turns Euclidean data into a
physical Hilbert-space structure and a time-translation operator
\cite{OS1,OS2,OsterwalderSeiler}.  In finite volume the relevant algebra is
elementary, but a formal development still has to identify which positive form
is physical, prove that the operator is forced by that form, and connect its
normalised spectrum to a mixing estimate whose constants do not deteriorate
with the spatial size.

This paper proves precisely that interface for the coupled $\mathbb Z_2$ slice
model already developed in the repository.  It does \emph{not} claim a new
Dobrushin estimate, a continuum limit, a Wightman reconstruction, or the first
formal treatment of reflection positivity.  In particular, recent Lean work
verifies Euclidean axioms for the free bosonic field
\cite{DouglasFormalQFT}.  Our narrower contribution is a concrete reconstructed
finite-volume transfer operator whose normalised fluctuation operator inherits
a single, explicit, volume-uniform gap exponent.

\section{Finite slice data}

For $L\in\N$, let
\[
  \SigmaL=(\operatorname{Fin} L\to\operatorname{Fin}2)
\]
be the finite set of spin configurations on a spatial slice.  Fix a positive
weight $w:\SigmaL\to\R_{>0}$ and the real symmetric spatial kernel
$K_\beta(\sigma,\tau)$.  The reconstructed site form on complex boundary
vectors is
\begin{equation}\label{eq:site-form}
  \langle u,v\rangle_{\mathrm{site},w}
   =\sum_{\sigma\in\SigmaL}
      \overline{u(\sigma)}v(\sigma)\,w(\sigma)^{-1}.
\end{equation}
The transfer operator forced by the site--bond identity is
\begin{equation}\label{eq:transfer}
  (T_{w,\beta}v)(\sigma)
   =w(\sigma)\sum_{\tau\in\SigmaL}K_\beta(\sigma,\tau)v(\tau).
\end{equation}
The symmetric weighted kernel is
\begin{equation}\label{eq:symmetric}
  S_{w,\beta}(\sigma,\tau)
  =\sqrt{w(\sigma)}K_\beta(\sigma,\tau)\sqrt{w(\tau)}.
\end{equation}
All three objects precede the present module.  The missing step was to state
and verify the exact real-to-complex dressing through which the Dobrushin lane
can be read on the reconstructed operator.

Define
\begin{equation}\label{eq:qembed}
  (Q_w u)(\sigma)=\sqrt{w(\sigma)}\,u(\sigma),
  \qquad u:\SigmaL\to\R,
\end{equation}
where the result is regarded as a complex boundary vector.  Since $w>0$, this
is the real restriction of the already available complex linear equivalence
given by multiplication by $\sqrt w$.

\section{The six machine-checked interfaces}

The new file \texttt{YangMills/OS/OSReconstructionUniform.lean} contains one
definition and the following six theorem declarations.

\begin{center}
\small
\begin{tabular}{@{}P{0.34\linewidth}P{0.58\linewidth}@{}}
\toprule
Lean declaration & Mathematical content \\
\midrule
\texttt{siteForm\_qEmbed} & $Q_w$ is an isometry from the real Euclidean pairing to \eqref{eq:site-form}. \\
\texttt{transferOp\_qEmbed} & $T_{w,\beta}Q_w=Q_wS_{w,\beta}$ on real boundary vectors. \\
\texttt{transferOp\_iterate\_qEmbed} & The same conjugation holds for every finite iterate. \\
\texttt{transferOp\_qEmbed\_tilt} & Perron normalisation corresponds to the Dobrushin tilt in one step. \\
\texttt{transferOp\_qEmbed\_tilt\_iterate} & The normalised conjugation holds for every finite iterate. \\
\texttt{os\_reconstruction\_uniform\_gap} & One positive exponent and the full reconstructed identity hold for all $L$. \\
\bottomrule
\end{tabular}
\end{center}

The first identity is exact:
\begin{proposition}[Dressing isometry]
For $w(\sigma)>0$ and real $u,v$,
\[
  \langle Q_wu,Q_wv\rangle_{\mathrm{site},w}
   =\sum_{\sigma\in\SigmaL}u(\sigma)v(\sigma).
\]
\end{proposition}
Indeed, each summand contains
$\sqrt{w(\sigma)}^2/w(\sigma)=1$.  In Lean the positivity hypothesis supplies
both the square-root identity and the nonzero denominator.

The second identity is the actual bridge:
\begin{proposition}[Forced-operator conjugation]
For every real boundary vector $u$,
\begin{equation}\label{eq:conjugation}
  T_{w,\beta}(Q_wu)=Q_w(S_{w,\beta}u).
\end{equation}
Consequently, for every $n\in\N$,
\[
  T_{w,\beta}^{n}(Q_wu)=Q_w(S_{w,\beta}^{n}u).
\]
\end{proposition}
The one-step proof reuses the previously checked complex conjugation theorem;
the iterate theorem is induction on $n$.  Thus the new file introduces no
parallel definition of the physical transfer operator.

\section{Perron tilt and the uniform endpoint}

Let $\lambda\in\R$ be the Perron eigenvalue and define the real tilt matrix
\begin{equation}\label{eq:tilt}
  P_{w,\beta,\lambda}(\sigma,\tau)
   =\lambda^{-1}S_{w,\beta}(\sigma,\tau).
\end{equation}
Scalar multiplication on the complex function space and matrix action on the
real function space then satisfy
\begin{equation}\label{eq:tilt-conj}
  \lambda^{-1}T_{w,\beta}(Q_wu)
   =Q_w(P_{w,\beta,\lambda}u),
\end{equation}
and the same equation holds after every finite iterate.  Formally,
\eqref{eq:tilt-conj} is division-free as a statement about the inverse scalar:
the proof merely distributes the same inverse through a finite sum.  The final
endpoint supplies $\lambda>0$.

For the coupled-slice weight $w=\operatorname{sliceW}(\gamma,L)$, the existing
Dobrushin endpoint produces $m>0$ independent of $L$, a positive Perron value
$\lambda_L$, and a positive vector $\Omega_L$.  The new composition is the
following statement.  The library parameter $L$ indexes the configuration
space $\Sigma_{L+1}$ in this endpoint; we preserve that convention instead of
silently rebasing the formal binder.

\begin{theorem}[Volume-uniform reconstructed gap]
Let $0<\alpha<1$ and suppose
\[
  2\tanh|\beta|+2\tanh|\gamma|\leq\alpha.
\]
Then there is $m>0$ such that for every $L\in\N$ there are
$\lambda_L>0$ and $\Omega_L:\Sigma_{L+1}\to\R_{>0}$ with
\begin{align*}
  \sum_\tau P_L(\sigma,\tau)\Omega_L(\tau)&=\Omega_L(\sigma),\\
  \bigl\|\operatorname{projectedTransfer}(\operatorname{opOf}P_L,
       \operatorname{vacOf}\Omega_L)\bigr\|&\leq e^{-m},
\end{align*}
and, for every real boundary vector $u$,
\[
  \lambda_L^{-1}T_{\operatorname{sliceW}(\gamma,L),\beta}
       (Q_{\operatorname{sliceW}(\gamma,L)}u)
  =Q_{\operatorname{sliceW}(\gamma,L)}(P_Lu).
\]
Here $P_L=P_{\operatorname{sliceW}(\gamma,L),\beta,\lambda_L}$.
\end{theorem}

The quantifier order is essential: $m$ is chosen before $L$.  This is the
precise sense in which the reconstructed gap is volume-uniform.

\section{Relation to prior work}

The classical reconstruction and reflection-positivity literature supplies
the conceptual framework \cite{OS1,OS2,OsterwalderSeiler,FILS}.  Dobrushin's
uniqueness theory supplies the high-temperature contraction mechanism
\cite{Dobrushin1968}.  Neither fact is claimed as new here.

\begin{center}
\small
\begin{tabular}{@{}P{0.25\linewidth}P{0.29\linewidth}P{0.36\linewidth}@{}}
\toprule
Work & Machine-checked content & Difference from the present endpoint \\
\midrule
Douglas et al.~\cite{DouglasFormalQFT} & Euclidean/Glimm--Jaffe axioms for a free bosonic field & Axiom verification, not a concrete reconstructed transfer operator with a measured uniform gap. \\
Classical OS papers \cite{OS1,OS2} & Analytic reconstruction theorem & Not proof-assistant checked and not specialised to this finite coupled-slice kernel. \\
Osterwalder--Seiler \cite{OsterwalderSeiler} & Lattice gauge reconstruction and positivity & Classical framework; no Lean object or repository-wide axiom report. \\
This paper & Six Lean interfaces joining a forced reconstructed operator to an existing Dobrushin gap & Finite $\mathbb Z_2$ slices only; no continuum or Wightman limit. \\
\bottomrule
\end{tabular}
\end{center}

A survey of public formal artifacts performed on 4 August 2026 found no prior
artifact with this exact conjunction: a concrete OS-reconstructed transfer
operator, an explicit Dobrushin window, and a gap exponent quantified before
the spatial volume.  Because repository-level and unindexed work can be missed,
we use this only as a qualified comparison, never as an unqualified ``first OS
reconstruction'' claim.

\section{Verification and reproducibility}

The authoritative source is the Lean file, not this prose.  Verification used
a Colab Pro+ CPU, high-RAM runtime with GPU disabled.  The clean target build
completed 8191 jobs with exit code zero between 16:25:28Z and 16:25:39Z on
4 August 2026.  The source contained exactly six theorem declarations.  Its
SHA-256 digest was
\[
  \ModuleHash,
\]
and the materialised \texttt{.olean} digest was
\[
  \OleanHash.
\]
The repository-wide oracle count for the frozen object is
\OracleReports.  A separate manifest records the branch SHA, byte counts,
line-ending hashes, capture method, full command transcript, and the attacks
attempted.  Successful compilation is evidence of elaboration, not a terminal
audit verdict; the frozen object is intended for independent review.

\section{Limitations}

The result is finite-volume and model-specific.  It does not construct a
continuum field theory, prove the full OS reconstruction theorem in generality,
or establish a Yang--Mills mass gap.  The physical space used here is the
finite reconstructed site space already identified by the reflection-positive
slice construction.  The analytic small-coupling window and the operator-norm
estimate are inherited without alteration from the Dobrushin lane.  Finally,
the formalisation does not turn a bibliographic survey into a mathematical
priority theorem; comparison language remains deliberately qualified.

\section{Conclusion}

The square-root dressing $Q_w$ is the missing typed interface between the
reflection-positive reconstruction and the symmetric kernel used by the
Dobrushin argument.  Once that interface is stated, the one-step equality,
its iterates, the Perron tilt, and the volume-uniform endpoint form a short but
fully checkable chain.  The resulting theorem says exactly what is uniform,
exactly which reconstructed operator carries the gap, and exactly which parts
are inherited.

\begin{thebibliography}{9}

\bibitem{OS1}
K.~Osterwalder and R.~Schrader,
\newblock Axioms for Euclidean Green's functions,
\newblock \emph{Communications in Mathematical Physics} \textbf{31} (1973), 83--112.
\newblock \href{https://doi.org/10.1007/BF01645738}{doi:10.1007/BF01645738}.

\bibitem{OS2}
K.~Osterwalder and R.~Schrader,
\newblock Axioms for Euclidean Green's functions II,
\newblock \emph{Communications in Mathematical Physics} \textbf{42} (1975), 281--305.
\newblock \href{https://doi.org/10.1007/BF01608978}{doi:10.1007/BF01608978}.

\bibitem{OsterwalderSeiler}
K.~Osterwalder and E.~Seiler,
\newblock Gauge field theories on a lattice,
\newblock \emph{Annals of Physics} \textbf{110} (1978), 440--471.
\newblock \href{https://doi.org/10.1016/0003-4916(78)90039-8}{doi:10.1016/0003-4916(78)90039-8}.

\bibitem{FILS}
J.~Fr\"ohlich, R.~Israel, E.~H.~Lieb, and B.~Simon,
\newblock Phase transitions and reflection positivity. I. General theory and long range lattice models,
\newblock \emph{Communications in Mathematical Physics} \textbf{62} (1978), 1--34.
\newblock \href{https://doi.org/10.1007/BF01940327}{doi:10.1007/BF01940327}.

\bibitem{Dobrushin1968}
R.~L.~Dobrushin,
\newblock The description of a random field by means of conditional probabilities and conditions of its regularity,
\newblock \emph{Theory of Probability and its Applications} \textbf{13} (1968), 197--224.
\newblock \href{https://doi.org/10.1137/1113026}{doi:10.1137/1113026}.

\bibitem{DouglasFormalQFT}
M.~R.~Douglas, S.~Hoback, A.~Mei, and R.~Nissim,
\newblock Formalization of QFT,
\newblock \href{https://arxiv.org/abs/2603.15770}{arXiv:2603.15770} (2026).

\end{thebibliography}

\end{document}
